% ============================================================================
%  Formulario T-19. Cartera de innovaciones
%  audIT Soluciones Tecnológicas SpA - Licitación TFEP-01/2026
% ============================================================================

\begin{formulario}{T-19}

\fichaInnovacion{
  tipo = {1},
  nombre = {Portal del transportista con liquidación en curso y expediente exportable verificable},
  problema = {El 60{,}4\,\% de la flota (226 de 374 tractocamiones) y el 56{,}8\,\% de los conductores (258 de 454) pertenecen a 148 transportistas subcontratados independientes (Caso, numeral 2.2). El 11\,\% de las liquidaciones mensuales se corrige con posterioridad a su emisión y el ciclo de cálculo demora 9 días tras el cierre de mes (Caso, numeral 7.3, p.~15). Los transportistas desconfían del equipamiento de monitoreo al percibirlo como una fiscalización punitiva unilateral si no les reporta un beneficio probatorio directo (Caso, p.~17).},
  tecnologia = {Portal web responsivo con aislamiento criptográfico por tenant e identidad federada (Microsoft Entra ID / OAuth 2.0). Emisión de expediente digital firmado electrónicamente con estampado de tiempo RFC 3161 y mecanismo de verificación pública en línea mediante código QR y URI segura sin exponer credenciales internas ni datos comerciales sensibles.},
  madurez = {Nivel TRL 9 para la firma electrónica avanzada y estampados cronológicos verificables en producción comercial en Chile bajo la Ley 19.799 \parencite{ley19799}. Nivel TRL 8 para la emisión de credenciales verificables descentralizadas conforme al estándar internacional W3C Verifiable Credentials Data Model v2.0 \parencite{w3c_vcdm2}.},
  fuentes = {W3C (2025). Verifiable Credentials Data Model v2.0; Biblioteca del Congreso Nacional de Chile (2002). Ley N.º 19.799 sobre Documentos Electrónicos, Firma Electrónica y Servicios de Certificación.},
  arquitectura = {Se inserta transversalmente entre la Capa 1 (Presentación: Portal Transportista), la Capa 4 (Servicios de Negocio: Servicio de Liquidación y Expediente Digital) y la Capa 7 (Seguridad: Azure Key Vault con módulo criptográfico HSM FIPS 140-2 Nivel 3 para sellado digital de evidencias).},
  edt = {Paquetes EDT 3.5 (Portal Web de Clientes y Transportistas) y EDT 5.3 (Módulo de Liquidación en Curso y Expediente Digital).},
  mes = {Mes 9 (prototipo funcional en laboratorio) y Mes 14 (habilitación en marcha blanca).},
  inversion = {Desarrollo de módulos de visualización en tiempo real, pasarela de verificación pública y suscripción de certificados corporativos de firma electrónica avanzada.},
  costooperacional = {Mantenimiento de infraestructura web, renovación anual de certificados criptográficos y soporte técnico a usuarios externos.},
  beneficio = {Reducción de liquidaciones corregidas del 11\,\% actual a menos del 2\,\%; reducción del ciclo de liquidación de 9 días a disponibilidad continua viaje a viaje; y adopción activa por al menos el 40\,\% de los transportistas subcontratados durante la marcha blanca.},
  indicador = {Línea base: 11\,\% de liquidaciones corregidas y 0 expedientes digitales. Meta: $\le 2\,\%$ de liquidaciones rectificadas y $\ge 40\,\%$ de transportistas emitiendo al menos un expediente probatorio al mes 15.},
  medicion = {Tercer cierre mensual posterior al paso a producción en régimen normal.},
  riesgo = {Desinterés inicial de los transportistas o baja comprensión de la utilidad probatoria del expediente digital (Probabilidad media, Impacto medio).},
  mitigacion = {Taller de inducción práctico en terminal durante el enrolamiento de flota, asistiendo la primera descarga de liquidación y expediente.},
  contingencia = {En caso de baja adopción en portal, el sistema despacha automáticamente el resumen mensual en PDF certificado vía correo electrónico y mensajería segura.}
}

\fichaInnovacion{
  tipo = {2},
  nombre = {Gemelo Digital operacional para simulación logística, ruteo predictivo y mantenimiento de flota},
  problema = {Asimetría y dispersión del 19\,\% en el rendimiento de combustible entre tractocamiones del mismo modelo operando en rutas idénticas (Caso, numeral 7.3, p.~15), sumado a la falta de anticipación operativa ante cierres de pasos cordilleranos que alcanzan hasta 12 días continuos (RT-10.05). La torre asigna fletes mediante heurísticas estáticas sin proyectar en tiempo real el consumo térmico ni la degradación mecánica.},
  tecnologia = {Plataforma Azure Digital Twins estructurada bajo el lenguaje Digital Twins Definition Language (DTDL v2), articulada con modelos cinemáticos vehiculares y telemetría de motor CAN/FMS (SAE J1939). Ingesta desacoplada mediante Azure IoT Hub y Apache Kafka, correlacionando perfil topográfico de elevación, gradiente térmico de montaña y carga transportada.},
  madurez = {Nivel TRL 7 (sistema integrado y validado en entorno operativo). Sustentado en el estándar DTDL del Digital Twin Consortium y en la norma ISO 23247 para gemelos digitales en sistemas de transporte y manufactura.},
  fuentes = {Digital Twin Consortium (2023). Digital Twin System Interoperability Framework; ISO (2021). ISO 23247 Automation systems and integration -- Digital twin framework for manufacturing.},
  arquitectura = {Se inserta transversalmente entre la Capa 4 (Servicios de Negocio: Planificación, Despacho y Mantenimiento Predictivo), la Capa 5 (Integración y Eventos: bus Kafka) y la Capa 6 (Datos: Lakehouse analítico Delta Lake y series temporales TimescaleDB).},
  edt = {Paquetes EDT 3.6 (Modelado DTDL de Flota y Rutas) y EDT 4.6 (Integración de Simulación en Torre de Control).},
  mes = {Mes 10 (piloto con 61 tractocamiones propios) y Mes 14 (integración en marcha blanca).},
  inversion = {Suscripción a servicios gestionados Azure Digital Twins, desarrollo de esquemas DTDL v2 y calibración de modelos predictivos de consumo.},
  costooperacional = {Consumo de operaciones de grafo y cómputo de simulación en Azure; compensado con ahorros en consumo de combustible y mantenimiento preventivo.},
  beneficio = {Reducción de la dispersión de combustible entre unidades idénticas del 19\,\% al 10\,\%; disminución del 15\,\% en fallas mecánicas en ruta; y anticipación de 6 horas en la reprogramación de convoyes ante contingencias climáticas.},
  indicador = {Línea base: 19\,\% de dispersión de rendimiento y cero anticipación predictiva. Meta: dispersión $\le 10\,\%$ y alertas predictivas de falla con al menos 48 horas de anticipación.},
  medicion = {Evaluación trimestral comparativa en régimen operativo sobre los 41 millones de km anuales.},
  riesgo = {Inconsistencias en telemetría de motor CANbus en tractocamiones subcontratados con interfaces dispares (Probabilidad media, Impacto medio).},
  mitigacion = {Entrenamiento del modelo predictivo acotado a los 61 tractocamiones propios con puerto FMS de fábrica activo antes de homologar terceros.},
  contingencia = {Si un vehículo pierde telemetría FMS, el gemelo digital conmuta a perfiles cinemáticos basados exclusivamente en coordenadas GPS y velocidad comercial.}
}

\fichaInnovacion{
  tipo = {3},
  nombre = {Operación desconectada resiliente con almacenamiento a bordo de 8 GB y vista única de flota},
  problema = {41 millones de kilómetros anuales con rutas de más de 80 kilómetros en zonas de silencio celular absoluto (Caso, numeral 2.2, p.~6), 34 camiones sin equipo y 340 unidades dispersas en tres plataformas de posicionamiento incompatibles, una sin interfaz de exportación (p.~34). Esta dispersión produce ceguera en la torre y riesgo crítico de multas laborales por pérdida de evidencia de jornada.},
  tecnologia = {Arquitectura híbrida Edge-to-Cloud con almacenamiento no volátil flash industrial $\ge 8$~GB en cabina, motor embebido SQLite 3 en modo WAL (\textit{Write-Ahead Logging}), compresión binaria Protocol Buffers y bus asíncrono MQTT v5.0 / Kafka. Emisión en contingencia de Documentos Electrónicos de Transporte (DET) con folios CAF y Timbre Electrónico DTE (TED) precargados en memoria criptográfica local antes del movimiento vehicular. Normalización de plataformas externas vía conectores API/Webhooks sin tocar hardware ajeno (Restricción 3).},
  madurez = {Nivel TRL 8/9 conforme a ISO 16290 \parencite{iso16290}. Sustentado en estándares consolidados de comunicación vehicular (SAE J1939-71, MQTT v5.0) y modelos formales de replicación distribuida con consistencia eventual determinista \parencite{kleppmann2017}.},
  fuentes = {Kleppmann, M. (2017). Designing Data-Intensive Applications; OASIS (2019). MQTT Version 5.0 Standard; FMS Standard (2025). Technical Specification rFMS version 5.0.0.},
  arquitectura = {Se inserta en la Capa 1 (Borde Terrestre: flash 8 GB, SQLite WAL y acoplador inductivo CAN), Capa 2 (Borde y Exposición: API Gateway con mTLS), Capa 5 (Eventos: Kafka Event Hubs) y Capa 6 (Datos: TimescaleDB y deduplicación por clave de idempotencia en PostgreSQL 16).},
  edt = {Paquetes EDT 3.4 (Conector Ingestión Telemática), EDT 4.2 (Instalación de Buffer a Bordo 8 GB) y EDT 4.5 (Acopladores Inductivos FMS J1939).},
  mes = {Mes 8 (flota propia completa 148 unidades intervenidas) y Mes 12 (despliegue general).},
  inversion = {Provisión de 182 computadores embarcados con flash industrial 8 GB, acopladores inductivos CAN sin contacto y desarrollo de adaptadores para 3 plataformas.},
  costooperacional = {Tarjetas SIM celulares multicarrier con consumo adaptativo ($\approx 15$~MB/mes) y cargos por mensajes satelitales Iridium SBD en emergencias.},
  beneficio = {Cero pérdida de registros de telemetría y jornada durante aislamientos de hasta 288 horas (12 días continuos de cierre invernal en Los Libertadores según RT-10.05); y visibilidad unificada del 100\,\% de los 374 tractocamiones en la consola de la torre de control.},
  indicador = {Línea base: ceguera total en sombras celulares $> 2$ horas y 3 plataformas desconectadas. Meta: 0\,\% pérdida de registros tras 72 h offline y 100\,\% de flota visible en consola unificada.},
  medicion = {Prueba destructiva de corte de enlace en marcha blanca (mes 14) y auditoría mensual de completitud de datos de viaje.},
  riesgo = {Reticencia de un proveedor telemático de terceros a otorgar API de integración o bloqueo de exportación (Probabilidad media, Impacto alto).},
  mitigacion = {Exigencia contractual de estándar técnico de homologación en convenio de adhesión y captura mediante reporte periódico automatizado.},
  contingencia = {Si un camión de tercero carece de integración telemática, el conductor utiliza la app móvil de audIT como canal de posicionamiento de respaldo durante el flete.}
}

\fichaInnovacion{
  tipo = {4},
  nombre = {Esquema de adhesión tecnológica en comodato con soberanía de datos para el transportista},
  problema = {226 camiones subcontratados cuyos dueños temen que el equipamiento telemático sea un mecanismo de fiscalización invasiva o que delate operaciones comerciales con la competencia (Caso, p.~17). 34 unidades no tienen ningún dispositivo y los transportistas carecen de capital para financiar hardware industrial de 8 GB.},
  tecnologia = {Convenio contractual de adhesión tecnológica bajo figura de comodato con opción de transferencia de dominio al cumplimiento del contrato (56 meses). Gobernanza de datos basada en políticas de soberanía bajo la Ley 21.719 \parencite{ley21719}: el dispositivo sólo transmite coordenadas mientras el camión ejecuta una orden de transporte activa para Curimón S.A., cesando automáticamente fuera del flete.},
  madurez = {Nivel TRL 9 en modelos contractuales y marcos de privacidad de datos personales regulados por la Ley 21.719 y el Código del Trabajo.},
  fuentes = {Congreso Nacional de Chile (2024). Ley N.º 21.719 sobre Protección y Tratamiento de Datos Personales; Ministerio del Trabajo (2006). Ley N.º 20.123 sobre Subcontratación.},
  arquitectura = {Se inserta en la Capa 4 (Servicios de Negocio: Módulo de Contratos y Consentimientos) y la Capa 7 (Seguridad: motor de políticas de privacidad ABAC y bitácora forense WORM).},
  edt = {Paquetes EDT 1.4 (Estrategia de Adhesión y Convenios de Comodato) y EDT 2.3 (Programa de Capacitación y Enrolamiento de Terceros).},
  mes = {Mes 4 (lanzamiento del programa de adhesión) y Mes 8 (cierre de suscripciones iniciales).},
  inversion = {Redacción legal de convenios marco de comodato, material de difusión y gestión de entrega de equipos para las 34 unidades sin dispositivo.},
  costooperacional = {Costo de administración de contratos y auditoría jurídica semestral de consentimientos de privacidad.},
  beneficio = {Adhesión voluntaria del 100\,\% de las 34 unidades desprovistas de equipo; eliminación del riesgo de demandas laborales por intromisión ilegítima en la privacidad de 258 conductores externos.},
  indicador = {Línea base: 0 convenios de adhesión y 34 unidades ciegas. Meta: 100\,\% de las 34 unidades incorporadas al programa y 0 reclamos por vulneración de privacidad de datos.},
  medicion = {Auditoría de adhesión al cierre de la Etapa 1 (mes 12) y revisión semestral de conformidades.},
  riesgo = {Desconfianza de gremios de transportistas o negativa a firmar el convenio de comodato (Probabilidad media, Impacto alto).},
  mitigacion = {Mesa de diálogo inicial con las asociaciones de transportistas exhibiendo las cláusulas de cese automático de transmisión y entrega del expediente digital gratuito.},
  contingencia = {Si un transportista rechaza el comodato, se le autoriza a homologar un equipo propio siempre que cumpla el estándar de exportación de telemetría de audIT.}
}

\fichaInnovacion{
  tipo = {5},
  nombre = {Sistema de bienestar, alerta preventiva de fatiga y ruteo hacia paradores seguros},
  problema = {La fatiga y somnolencia al volante es la principal causa de accidentes graves en el transporte de carga interurbano (numeral 7.3 del Caso). Los conductores enfrentan dificultades para encontrar estacionamientos seguros con servicios básicos en la red vial chilena, exponiéndose a robos o extendiendo la conducción más allá del límite legal de 5 horas continuas del Art.~25 bis.},
  tecnologia = {Algoritmo de ruteo preventivo y catálogo georreferenciado de Paradores Seguros certificados en las Rutas 5 Norte, 5 Sur y 68. Integración de tiempos de marcha acumulados en el dispositivo a bordo con alerta sonora y visual en cabina emitida con 45 minutos de anticipación al límite de descanso, indicando la ubicación del parador seguro más próximo con plazas disponibles.},
  madurez = {Nivel TRL 8, sustentado en algoritmos geoespaciales consolidados de optimización de rutas (Open Source Routing Machine / PostGIS) y catálogos de infraestructura vial del MOP.},
  fuentes = {Ministerio de Obras Públicas (2024). Catálogo Nacional de Áreas de Descanso y Servicios Viales; Dirección del Trabajo (2009). Resolución Exenta N.º 1213 sobre control de asistencia y descansos en transporte interurbano.},
  arquitectura = {Se inserta en la Capa 1 (Borde Terrestre: display de cabina con alertas de descanso), Capa 4 (Servicios de Negocio: Servicio de Jornada y Ruteo Vial) y Capa 6 (Datos: base geoespacial PostGIS de paradores seguros).},
  edt = {Paquetes EDT 3.7 (Módulo de Gestión de Jornada y Paradores Seguros) y EDT 4.7 (Piloto de Alertas en Cabina).},
  mes = {Mes 6 (catálogo de paradores) y Mes 12 (activación de alertas en flota activa).},
  inversion = {Levantamiento en terreno y certificación de paradores seguros en las rutas principales, desarrollo del motor de recomendación en cabina y app móvil.},
  costooperacional = {Mantenimiento y actualización mensual de la base de datos de paradores y servicios viales habilitados.},
  beneficio = {Reducción del 30\,\% en infracciones por exceso de jornada continua (Art.~25 bis); cero accidentes por volcamiento o colisión atribuibles a fatiga en la flota intervenida; y elevación del índice de satisfacción del conductor.},
  indicador = {Línea base: 14 infracciones anuales de jornada y cero alertas preventivas en cabina. Meta: 0 infracciones de jornada continua y $\ge 95\,\%$ de descansos efectuados en paradores seguros certificados.},
  medicion = {Monitoreo mensual de eventos de jornada y reporte semestral de siniestralidad vial en comités paritarios.},
  riesgo = {Falta de plazas disponibles en los paradores recomendados en horario nocturno (Probabilidad media, Impacto medio).},
  mitigacion = {Convenios preferenciales de reserva de estacionamiento con concesionarias de autopistas y servicentros de las rutas críticas.},
  contingencia = {Si un parador se encuentra colapsado, el sistema recalcula en tiempo real la ruta hacia la siguiente instalación segura autorizada sin penalizar al conductor.}
}

\end{formulario}
